\documentclass[../main]{subfiles}
\begin{document}
\chapter{Intensidad de campo magnético}
\info{
Contenido}{
Intensidad de campo, inducción, permeabilidad magnética absoluta, del vacío y relativa.
}
\img{images/02/bobinas}{Forma habitual de una bobina.}
\info {Intensidad de campo o excitación magnética en un solenoide o bobina}{
La Intensidad de campo magnético $H$ es la relación que existe entre el producto del número de vueltas $N$ de una bobina, por la intensidad de la corriente que circula $I$ dividido la longitud del núcleo de dicha bobina $l$.
\begin{equation*}
    H=\frac{N \times I}{l}
\end{equation*}

    siendo:
    
    $H$: Intensidad de campo magnético en ampere vuelta sobre metro $[\frac{Av}{m}]$.
    
    $N:$ Número de vueltas o número de espiras.
    
    $l:$ Longitud del solenoide o bobina en metros. $[m]$
}
\info{Induccción Magnética}{
\begin{equation*}
    B=\mu \times H
\end{equation*}
siendo

$B$: Inducción magnética en Teslas $[T]$

$H$: Intensidad de campo magnético $[Av/m]$

$\mu$: Permeabilidad magnética $[\frac{T \cdot m}{A}]$.


\begin{equation*}
    \mu_0= 4\times \pi \times 10^{-7} [\frac{T\cdot m}{A}]
\end{equation*}

$\mu_0$: Permeabilidad magnética del aire o vacío,

}

\info{Permeabilidad relativa}{
Para comparar entre sí los materiales, se entiende la permeabilidad magnética absoluta $\mu$ como el producto entre la permeabilidad magnética relativa $\mu _{r}$ y la permeabilidad magnética de vacío $\mu _{0}$

\begin{equation*}
    \mu_=\mu_r \times \mu_0
\end{equation*}

siendo:

$\mu_r$: Permeabilidad relativa (es adimensional)

$\mu_0$: Permeabilidad del vacío.$[\frac{T\cdot m}{A}]$

$\mu$: Permeabilidad absoluta. $[\frac{T\cdot m}{A}]$

}


\actividad{Resolver los siguientes problemas}
\begin{enumerate}
    \item Un campo magnético uniforme tiene en el aire $\mu_0=4\pi \times 10^{-7} \frac{T \cdot m}{A}$ una inducción $B=12000Gauss$. Calcular la intensidad de campo magnético $H$.
    \item Un solenoide de $N=400$ espiras y longitud $l=500cm$ está recorrido por una corriente eléctrica de $I=10A$. Si el núcleo es de aire. Calcular la intensidad de campo magnético $H$ desde el interior del solenoide.
    \item Un solenoide de $N=200$ espiras y longitud de $l=40cm$ está bobinado sobre un núcleo de madera de $r=3cm$ de radio. La intensidad de la corriente por el solenoide es de 10A. Calcular:
    \begin{enumerate}
        \item Intensidad del campo magnético $H$ en el interior de solenoide.
        \item La inducción magnética B en el núcleo del solenoide.
        \item Flujo magnético en el núcleo $\Phi$.
    \end{enumerate}
    \item La intensidad de campo magnético en el interior del solenoide es de $H=2000\frac{Av}{m}$ y su longitud es de $l=30cm$. Calcular:
    \begin{enumerate}
        \item Amperio vuelta del solenoide.
        \item Intensidad de corriente $I$, en el conductor del solenoide si tiene $N=1000$ espiras.
    \end{enumerate}
    \item Calcular la intensidad de corriente que debe circular por un solenoide de $N=500$ espiras y longitud de $l=40cm$, para que la intensidad de campo en el núcleo sea de $H=4000\frac{Av}{m}$.
    %\item La intensidad de campo magnético es de $H=12\frac{Av}{m}$ y la permeabilidad relativa $\mu_r=$    ... en este material ... FALTA
    \item Calcular la intensidad de campo magnético $H$ en el interior un núcleo de acero que con una inducción magnética de $B=0,6T$ tiene una permeabilidad relativa de $\mu_r=6000$
\end{enumerate}
\end{document}